\section{问题二：第二检测点选择}

\subsection{问题分析}

问题一说明了如何由观测确定可能位置，但一次测向形成的角域通常仍较宽。要继续缩小区域，机器狗需前往第二个检测点。选点过近会使两条示向线近乎平行；只追求较大交会角，又可能走得过远或失去信号。因此，本问先确定全向源再次接收的候选区域，再比较其中各点的交会质量和移动代价。

\subsection{模型建立：总体框架}

选址分为“确定可行范围”和“在范围内择优”两步。首先，利用首次示向及成功接收这一事实，求与观测相容的源位置和接收半径，构造第二检测点的保证接收区。其次，在该区域内比较最坏交会角、定位误差代理、接收裕量和移动距离。图~\ref{fig:p2-workflow} 展示这一顺序；筛选先于评分，是为了避免只因交会角大就选中可能失联的检测点。数值实现使用离散采样，故下文算例的保证接收结论限于采样范围。

% 制图备注（problem2_workflow.pdf；仅供后续重绘，不改当前图片）：
% 目的：突出“两层决策”而不只是并列罗列五个符号，即先保证可接收、后比较定位质量。
% 第一层输入框写 S_1、alpha±epsilon、目标圆 T、成功接收、R_G∈[1000,1500] m；
% 由此指向源位置可行域 Omega_1，旁注 20 m 内可直接光学定位的分支，不进入二次选点。
% 对 Omega_1 的每个 G 计算 rho(G)=max{1000,|G-S_1|}，再交所有 B(G,rho(G)) 得 C_safe；
% 画“候选网格 + 最短基线 100 m”硬筛选门，明确 M(s)>=0 才属于保证接收模式。
% 若保证接收网格为空，单独画虚线退化分支 C_poss，并标“可能接收、须报告风险”，不可与主路径混为一谈。
% 第二层只对保留点计算 gamma_min、E_max、M、L，先标量纲归一化，再由 J 排序；
% 最后一框沿首次示向线法向分组，输出 S_2^(+) 与 S_2^(-)，并给出各点模式与指标。
% 箭头上注明“源位置样本”“候选点”“最坏情形指标”等数据流；主路径用实线、退化分支用虚线，中文说明应可单栏阅读。
\safeimage[0.98\textwidth]{figures/problem2_workflow.pdf}{问题二保证接收优先的鲁棒选点流程\label{fig:p2-workflow}}

\subsubsection{候选区域构造}

\paragraph{步骤一：刻画首次测向后的源位置}

源位置未知，首次示向又有 \(1\degree\) 的误差，故第二检测点不能只针对一条中心射线设计。本文借鉴最坏情形方位传感器布设思想\cite{tokekar2013bearing}，结合题设的目标圆、有效接收半径、\(20\,\mathrm{m}\) 光学定位距离和机器狗移动成本，构造可计算的选址区域。

首先用题设目标圆限定空间搜索范围，记为
\begin{equation}
  \mathcal T=\{g\in\R^2:\|g\|_2\le1800\}.
\end{equation}
第一次检测点为 \(S_1\)，测得示向度为 \(\alpha\)，误差上界为 \(\varepsilon=1\degree\)。用 \(\operatorname{wrap}\) 将角度差映射到 \([-180\degree,180\degree)\)，则首次示向扇形为
\begin{equation}
  W_1=\left\{
  g:\left|\operatorname{wrap}\!\left(
  \arg(g-S_1)-\alpha
  \right)\right|\le\varepsilon
  \right\}.
\end{equation}
目标圆和测向扇形分别限定全局位置与方向，第一次成功接收还提供了距离信息：源到 \(S_1\) 的距离不超过实际有效接收半径，因而不超过 \(1500\,\mathrm{m}\)。若源距 \(S_1\) 不超过 \(20\,\mathrm{m}\)，机器狗可直接进行光学精确定位，不必选择第二检测点。排除这种情形后，待选址的源位置可行域为
\begin{equation}
  \boxed{
  \Omega_1=\mathcal T\cap W_1\cap
  \left\{g:20<\|g-S_1\|_2\le1500\right\}.
  }
  \label{eq:p2-source-region}
\end{equation}

式~\eqref{eq:p2-source-region} 保留与首次观测相容的全部源位置。为在程序中评价最坏情形，按方向和距离对该连续区域作确定性采样。对任一方向
\(\theta\in[\alpha-\varepsilon,\alpha+\varepsilon]\)，令
\(\vect u=\vect u(\theta)\)。射线 \(S_1+r\vect u\) 到目标圆边界的正向距离由
\(\|S_1+r\vect u\|_2=1800\) 得到：
\begin{equation}
  r_{\mathcal T}(\theta)=
  -S_1^{\mathsf T}\vect u+
  \sqrt{(S_1^{\mathsf T}\vect u)^2+
  1800^2-\|S_1\|_2^2}.
  \label{eq:p2-ray-circle}
\end{equation}
故该方向的径向范围为
\begin{equation}
  20<r\le
  r_{\max}(\theta)
  =\min\{1500,r_{\mathcal T}(\theta)\}.
\end{equation}
在每个采样方向上，径向距离从 \(20\,\mathrm{m}\) 取至 \(r_{\max}(\theta)\)，得到 \(\Omega_1\) 的离散表示。后续的接收约束和定位质量均在同一组源位置样本上计算，使各指标基于一致的不确定性范围。

\paragraph{步骤二：筛选保证接收的第二检测点}

记干扰源实际有效接收半径为 \(R_G\in[1000,1500]\)。若真实位置为 \(G\in\Omega_1\)，由第一次成功接收可知
\[
  R_G\ge\|G-S_1\|_2.
\]
若仅以题设的 \(1000\,\mathrm{m}\) 下界选点，就浪费了首次成功接收提供的信息。对给定源位置 \(G\)，实际接收半径还必须满足 \(R_G\ge\|G-S_1\|_2\)，故与首次观测相容的最小半径为
\begin{equation}
  \rho(G)=\max\{1000,\|G-S_1\|_2\}.
  \label{eq:p2-compatible-radius}
\end{equation}
若第二检测点 \(S_2\) 对每个 \(G\in\Omega_1\) 均满足 \(\|S_2-G\|_2\le\rho(G)\)，则所有与首次观测相容的源位置和接收半径都会在 \(S_2\) 处产生可接收信号。由此定义保证接收区
\begin{equation}
  \boxed{
  \mathcal{C}_{\mathrm{safe}}
  =\mathcal T\cap
  \bigcap_{G\in\Omega_1}
  B\!\left(G,\rho(G)\right).
  }
  \label{eq:p2-safe-region}
\end{equation}
式~\eqref{eq:p2-safe-region} 使用圆盘交集，是为了要求候选点对每一种可能源位置都能接收；只要违反其中一个圆盘约束，就存在第二次失联的相容情形。闭圆盘的交集仍为凸集。为避免两次测向基线过短，设
\(L_{\min}=100\,\mathrm{m}\)，首选候选区域为
\begin{equation}
  \mathcal{C}_0=
  \mathcal{C}_{\mathrm{safe}}\cap
  \{s:\|s-S_1\|_2\ge L_{\min}\}.
  \label{eq:p2-primary-candidates}
\end{equation}

当首次检测点靠近目标圆边界，或候选网格较粗时，离散的 \(\mathcal{C}_0\) 可能没有网格点。为使程序仍能返回备选方案，此时使用“可能接收区”
\begin{equation}
  \mathcal{C}_{\mathrm{poss}}
  =\mathcal T\cap
  \bigcup_{G\in\Omega_1}B(G,1500)
  \cap\{s:\|s-S_1\|_2\ge L_{\min}\},
  \label{eq:p2-possible-region}
\end{equation}
该区域中的点只需对至少一种相容源位置可接收，并不具备最坏情形保证。程序先搜索 \(\mathcal{C}_0\)；只有保证接收网格为空时，才转入 \(\mathcal{C}_{\mathrm{poss}}\)，并在结果中标记退化模式及接收风险。

\subsubsection{选点指标与综合决策}

\paragraph{步骤三：构造定位质量指标}

对给定第二检测点 \(s\) 和可能源位置 \(g\)，两条站源向量的锐交会角为
\begin{equation}
  \gamma(s,g)=
  \arccos
  \frac{\left|(g-S_1)^{\mathsf T}(g-s)\right|}
  {\|g-S_1\|_2\|g-s\|_2},
  \qquad 0\le\gamma\le90\degree.
  \label{eq:p2-crossing-angle}
\end{equation}
交会角越大，两条测向带的交会越稳定；接近零时，微小角度误差也可能造成较大的位置偏移。为避免只在某一个有利源位置上取得好结果，定义候选点的最坏交会角
\begin{equation}
  \gamma_{\min}(s)=
  \min_{g\in\Omega_1}\gamma(s,g).
\end{equation}

交会角并非唯一决定因素：相同的角度误差在更远距离处对应更大的横向偏差。令
\(d_1=\|g-S_1\|_2\)、\(d_2=\|g-s\|_2\)，两条测向带在目标附近的横向半宽近似为
\(d_i\tan\varepsilon\)。据此定义一阶保守误差代理
\begin{equation}
  E(s,g)=
  \frac{(d_1+d_2)\tan\varepsilon}
  {\sin\gamma(s,g)}.
  \label{eq:p2-error-proxy}
\end{equation}
式~\eqref{eq:p2-error-proxy} 同时惩罚长测距和小交会角；当 \(\gamma\) 趋近零时，误差代理迅速增大。若 \(s\) 距 \(g\) 不超过 \(20\,\mathrm{m}\)，可直接光学定位，程序将该情形的误差代理记为零。候选点在全部可能源位置中的最坏值定义为
\begin{equation}
  E_{\max}(s)=\max_{g\in\Omega_1}E(s,g).
\end{equation}

\paragraph{步骤四：综合接收、几何与移动代价}

候选点的最坏接收裕量定义为
\begin{equation}
  M(s)=\min_{g\in\Omega_1}
  \left[\rho(g)-\|s-g\|_2\right].
  \label{eq:p2-margin}
\end{equation}
\(M(s)\) 是候选点距最不利接收边界的有符号裕量；\(M(s)\ge0\) 等价于 \(s\in\mathcal{C}_{\mathrm{safe}}\)。移动距离记为 \(L(s)=\|s-S_1\|_2\)。四项指标的含义和优化方向见表~\ref{tab:p2-indicators}。

\begin{table}[H]
  \centering
  \caption{问题二候选点评价指标及其建模作用}
  \label{tab:p2-indicators}
  \small
  \begin{tabularx}{\textwidth}{cXcX}
    \toprule
    指标 & 物理含义 & 优化方向 & 模型作用 \\
    \midrule
    \(\gamma_{\min}\) & 最不利源位置下的锐交会角 & 增大 & 避免两条测向带近似平行 \\
    \(E_{\max}\) & 同时考虑距离与交会角的定位误差代理 & 减小 & 直接衡量最坏定位质量 \\
    \(M\) & 第二次检测的最坏接收裕量 & 增大 & 区分保证接收与失联风险 \\
    \(L\) & 从首次检测点到候选点的移动距离 & 减小 & 控制机器狗的时间代价 \\
    \bottomrule
  \end{tabularx}
\end{table}

四项指标的量纲不同。以 \(100\,\mathrm{m}\)、\(1000\,\mathrm{m}\) 和目标圆直径 \(3600\,\mathrm{m}\) 分别归一化误差代理、接收裕量和移动距离后，构造无量纲评分
\begin{equation}
  \begin{aligned}
  J(s)={}&
  \frac{0.65}{1+E_{\max}(s)/(100\,\mathrm{m})}
  +0.35\sin\gamma_{\min}(s)\\
  &+0.15\operatorname{clip}
  \left(\frac{M(s)}{1000\,\mathrm{m}},-1,1\right)
  -0.08\frac{L(s)}{3600\,\mathrm{m}}.
  \end{aligned}
  \label{eq:p2-score}
\end{equation}
误差代理权重最大，用于优先控制最不利源位置下的定位偏差；交会角补充几何质量，接收裕量鼓励远离失联边界，移动距离则在定位质量接近时起到择优作用。评分仅在同一候选模式内排序，不代替保证接收这一硬约束。在 \(\mathcal{C}_0\) 非空时，最优点为
\begin{equation}
  S_2^*=\arg\max_{s\in\mathcal{C}_0}J(s).
  \label{eq:p2-decision}
\end{equation}
由于首次示向线两侧都可能形成有效交会，程序分别输出两侧评分最高的候选点。在对称算例中二者等价；若实际场景中一侧受到圆域边界或障碍限制，可比较并选用另一侧。

\subsubsection{精确复核与计算边界}

式~\eqref{eq:p2-error-proxy} 是快速筛选候选点的一阶误差代理，不能直接解释为定位区域的真实直径。若要精确复核评分靠前的少量候选点，可进一步计算最坏情形下的定位区域直径：
\begin{equation}
  \begin{aligned}
  D_{\mathrm{wc}}(s)
  =\max_{g\in\Omega_1}
  \max_{e_2\in[-\varepsilon,\varepsilon]}
  \operatorname{diam}\bigl[
  &\mathcal T\cap W(S_1,\alpha)\\
  &{}\cap W(s,\arg(g-s)+e_2)
  \bigr].
  \end{aligned}
  \label{eq:p2-worst-diameter}
\end{equation}
式~\eqref{eq:p2-worst-diameter} 同时考虑真实源位置与第二次测向误差，但计算量较大，不适合直接用于整个候选网格。本次实现以 \(E_{\max}\) 排序，报告的数值均为一阶误差代理；\(D_{\mathrm{wc}}\) 仅作为后续精确复核的定义，不计入下文算例结果。

\subsection{模型求解：算法设计与复杂度}

算法~\ref{alg:p2} 给出离散实现。先生成与首次观测相容的源位置样本，再对目标圆内的候选网格执行接收筛选；随后计算保留点的最坏情形指标与评分，最后按示向线两侧输出推荐点。

\begin{modelalgorithm}{保证接收优先的第二检测点选择}
  \label{alg:p2}
  \AlgoIn{\(S_1\)、示向度 \(\alpha\)、误差 \(\varepsilon\)、圆域与接收半径参数、网格步长 \(h\)}
  \AlgoOut{候选模式、候选区域离散表示及示向线两侧的推荐点}
  \begin{enumerate}[leftmargin=2.5em]
    \AlgoStep{在 \([\alpha-\varepsilon,\alpha+\varepsilon]\) 上确定性采样方向，并按式~\eqref{eq:p2-ray-circle} 计算各方向径向上界。}
    \AlgoStep{在 \(20\,\mathrm{m}\) 至径向上界间采样，得到 \(\Omega_1\) 的离散点集。}
    \AlgoStep{在目标圆域内生成间距为 \(h\) 的候选网格，并删除基线短于 \(L_{\min}\) 的点。}
    \AlgoStep{对每个候选点遍历全部源位置样本，计算 \(M\)、\(\gamma_{\min}\)、\(E_{\max}\) 与 \(J\)。}
    \AlgoStep{若存在 \(M\ge0\) 的点，仅保留保证接收点；否则保留可能接收点并标记退化模式。}
    \AlgoStep{按第一次示向线的法向投影区分左右两侧，各返回一个评分最高的点。}
  \end{enumerate}
\end{modelalgorithm}

设源位置样本数为 \(N_\Omega\)，目标圆内的候选网格点数为 \(N_C\)。每个候选点都需遍历全部源位置样本，因此时间复杂度为 \(O(N_CN_\Omega)\)，空间复杂度为 \(O(N_C+N_\Omega)\)。减小网格步长 \(h\) 能细化候选边界与推荐点，但候选点数量及运行时间近似按 \(h^{-2}\) 增长。

\subsection{结果分析：算例结果与策略解释}

设第一次检测点为 \(S_1=(0,0)\,\mathrm{m}\)，示向度 \(\alpha=0\degree\)，即中心方向指向正东。取 \(9\) 个角度样本、\(31\) 个径向样本，候选网格步长为 \(100\,\mathrm{m}\)。程序生成 \(279\) 个源位置样本，并筛得 \(122\) 个保证接收候选点，其坐标的最小包围盒为
\[
  x\in[100,1000],\qquad y\in[-800,800].
\]
该矩形仅说明候选点的坐标范围，不代表矩形内每个位置都满足保证接收条件。离散候选点的形状及网格敏感性见下图。

% 制图备注（problem2_candidate_region.pdf；仅供后续重绘，不改当前图片）：
% 目的：左图回答“候选区域在哪里、推荐点为何选在两侧”，右图回答“网格细化后结论是否稳定”。
% 左图坐标单位 m，标 S_1=(0,0)、首次示向中心线 alpha=0°及 ±1° 边界；
% 以浅色显示 Omega_1 的 279 个源位置样本或其扇形包络，以另一种颜色显示 100 m 网格下的 122 个 C_safe 候选点；
% 可叠加由 M(s)=0 定义的接收安全边界，但须区分连续理论区域与离散采样近似；
% 用醒目标记 S_2^(+)=(700,700) m、S_2^(-)=(700,-700) m，并画出从 S_1 到两点的移动路径；
% 在推荐点旁注最坏交会角 41.1147°、误差代理 68.3232 m、接收裕量 23.8377 m；
% 图例必须解释源样本、保证接收候选点、首次示向线与推荐点，避免把源样本误读为检测点。
% 右图横轴为候选网格步长 200/100/50 m，建议同图展示最坏误差代理与综合评分的变化；
% 两指标量纲不同，使用分面或清晰标出的双轴，按表(\ref{tab:p2-sensitivity})的数值绘制而非凭示意连线；
% 标出三种步长下推荐点 (600,800)、(700,700)、(750,650) m，说明策略结构稳定而坐标略有细化。
% 主结论写在图注或图内简短说明：先满足采样意义下 M>=0，再选交会质量较好的双侧点；不用暗示连续域已获严格保证。
\safeimage[0.96\textwidth]{figures/problem2_candidate_region.pdf}{问题二保证接收候选区域、推荐点及网格敏感性}

示向线两侧的推荐点为
\begin{equation}
  \boxed{
  S_2^{(+)}=(700,700)\,\mathrm{m},\qquad
  S_2^{(-)}=(700,-700)\,\mathrm{m}.
  }
\end{equation}
目标圆、首次检测点及示向扇形均关于 \(x\) 轴对称，因此两个推荐点在本算例中的评分相同。实际执行时，可结合后续任务方向选择其中一侧。表~\ref{tab:p2-result} 列出单侧推荐点的评价结果。

\begin{table}[H]
  \centering
  \caption{问题二默认算例的推荐点质量}
  \label{tab:p2-result}
  \begin{tabular}{lc}
    \toprule
    指标 & 数值 \\
    \midrule
    候选模式 & 保证接收 \\
    移动距离 & \(989.9495\,\mathrm{m}\) \\
    纯移动时间 & \(197.99\,\mathrm{s}\) \\
    最坏交会角 & \(41.1147\degree\) \\
    最坏定位误差代理 & \(68.3232\,\mathrm{m}\) \\
    最坏接收裕量 & \(23.8377\,\mathrm{m}\) \\
    综合评分 & \(0.597888\) \\
    \bottomrule
  \end{tabular}
\end{table}

从 \(S_1\) 到任一推荐点需移动 \(989.9495\,\mathrm{m}\)。以题设速度 \(5\,\mathrm{m/s}\) 计算，纯移动时间为
\[
  T_{\mathrm{move}}=\frac{989.9495}{5}=197.99\,\mathrm{s}.
\]
到达后完成一次测向还需 \(5\,\mathrm{s}\)；若中途切换频道，则另计 \(1\,\mathrm{s}\)。最坏接收裕量为 \(23.8377\,\mathrm{m}>0\)，即在所用的 \(279\) 个离散源位置样本及其相容接收半径下，推荐点均保留正的接收余量。这是离散样本上的保证，连续可行域仍受采样密度限制。

推荐点的最坏交会角为 \(41.1147\degree\)，最坏定位误差代理为 \(68.3232\,\mathrm{m}\)，综合评分为 \(0.597888\)。这组结果表明，推荐位置并非单纯追求 \(90\degree\) 交会，而是在远近源位置的几何质量、再次接收条件和移动时间之间取得折中。这里的移动项权重服务于单次复测选址；完整搜索任务的时间分配须结合后续策略另行确定。

\subsubsection{网格敏感性与策略解释}

固定角度和径向采样，仅改变候选网格步长，得到表~\ref{tab:p2-sensitivity}。从 \(100\,\mathrm{m}\) 细化到 \(50\,\mathrm{m}\) 后，评分由 \(0.597888\) 增至 \(0.599272\)，约提高 \(0.23\%\)；误差代理由 \(68.3232\,\mathrm{m}\) 降至 \(66.7103\,\mathrm{m}\)，约降低 \(2.36\%\)。三种网格均选出“沿示向线前移并向侧方偏移”的位置，说明这一选址结构对网格步长较稳定，而推荐点坐标仍可随网格细化调整。

\begin{table}[H]
  \centering
  \caption{问题二候选网格的敏感性分析}
  \label{tab:p2-sensitivity}
  \small
  \begin{tabular}{cccccc}
    \toprule
    步长/m & 单侧推荐点/m & 最坏交会角/\(\degree\) &
    误差代理/m & 接收裕量/m & 评分 \\
    \midrule
    \(200\) & \((600,800)\) & \(41.5584\) & \(71.6096\) & \(11.5861\) & \(0.590467\) \\
    \(100\) & \((700,700)\) & \(41.1147\) & \(68.3232\) & \(23.8377\) & \(0.597888\) \\
    \(50\)  & \((750,650)\) & \(40.6668\) & \(66.7103\) & \(22.3199\) & \(0.599272\) \\
    \bottomrule
  \end{tabular}
\end{table}

只沿 \(S_1\) 处的示向垂线横向移动并非可靠规则：源距未知时，远端源仍可能与第二次示向形成很小的交会角。本算例推荐点同时向示向前方和侧方移动，兼顾了近端、远端与中间源位置。两点在离散样本上均满足非负接收裕量，因此优于仅凭交会角选择的经验位置。

% 可选补充图备注一：若需解释权重对结果的影响，可绘制选址性能地形图，而非替换现有候选区图。
% 横纵轴为第二检测点 x、y，底色表示同一可行模式内的 J(s)，M(s)=0 等值线标接收安全边界，
% E_max=70/80/100 m 画虚线等值线；用菱形标两推荐点，并在图例注明不可行区域不参与评分。
% 读图目标是识别推荐点如何在接收边界、误差代理和移动距离之间折中，避免把评分颜色画到失联点上。
%
% 可选补充图备注二：若需展示多目标权衡，绘制 Pareto 散点图：横轴移动时间 s，纵轴 E_max m，
% 颜色表示 M m，符号大小表示 gamma_min；只对 M>=0 的点绘制非支配前沿，另用灰点标不可保证接收者。
% 标出 200/100/50 m 网格的推荐方案及其评分，说明权重变化可能改变选点，但不改变保证接收优先的规则。

\subsection{结论}

第二检测点的选择将源位置的不确定性转化为可比较的行动方案。默认算例在 \(100\,\mathrm{m}\) 网格上推荐 \((700,\pm700)\,\mathrm{m}\)，两点在所用源位置样本中均满足再次接收条件。本问给出了比较单源复测点的基准；问题三、四还需同时安排未知频道扫描和多个已发现源的任务，其执行程序采用另行标定的复测点生成规则，未直接调用本问的保证接收筛选与评分。
